\documentclass[11pt,a4paper]{article}
\usepackage[utf8]{inputenc}
\usepackage[T1]{fontenc}
\usepackage[french]{babel}
\usepackage{amsmath,amsfonts,amssymb,geometry}
\geometry{hmargin=2cm,vmargin=2cm}

\newcommand{\N}{\mathbb{N}}
\newcommand{\Z}{\mathbb{Z}}
\newcommand{\Q}{\mathbb{Q}}
\newcommand{\R}{\mathbb{R}}
\newcommand{\D}{\mathbb{D}}
\newcommand{\C}{\mathbb{C}}

\newcounter{exocount}
\newenvironment{exercice}[1]{
    \stepcounter{exocount}
    \vspace{0.5cm}
    \noindent\textbf{Exercice \arabic{exocount} -- #1}
    \par\medskip
}{
    \par
}

\begin{document}

\noindent
\begin{minipage}[t]{0.5\textwidth}
    \textbf{LYCÉE POLYVALENT JACQUES BREL} \\
    Année scolaire 2026-2027
\end{minipage}
\begin{minipage}[t]{0.5\textwidth}
    \raggedleft
    \textbf{Niveau : Secondaire} \\
    \textbf{Date : } 25/10/2026
\end{minipage}

\vspace{0.2cm}
\noindent\rule{\textwidth}{0.5pt}
\vspace{0.3cm}

\begin{center}
    {\LARGE \bfseries Feuille d'exercices}
\end{center}

\vspace{0.1cm}
\noindent\rule{\textwidth}{1.5pt}
\vspace{0.4cm}

\noindent
\textbf{Nom : } \dots\dots\dots\dots\dots\dots\dots\dots\dots\dots\dots\dots\dots\dots\dots\quad 
\textbf{Prénom : } \dots\dots\dots\dots\dots\dots\dots\dots\dots\dots\dots\dots\dots\dots\dots

\vspace{0.4cm}
\begin{exercice}{Simplification de fraction (Seconde)}
Simplifier la fraction $\dfrac{42}{56}$.
\end{exercice}
\begin{exercice}{Appartenance aux ensembles (Seconde)}
Pour chacun des nombres, indiquer à quel(s) ensemble(s) $\N$, $\Z$, $\D$, $\Q$ et $\R$ il appartient.
\begin{itemize}
	\item $a = -5$
	\item $b = 0{,}4$
	\item $c = \sqrt{2}$
	\item $d = \frac{1}{3}$
	\item $e = \frac{1}{4}$
	\item $f = 0$
\end{itemize}
\end{exercice}
\begin{exercice}{Plus petit ensemble d'appartenance (Seconde)}
Pour chacun des nombres suivants, indiquer quel est le plus petit ensemble auquel il appartient.
\begin{itemize}
	\item $a = -\frac{7}{2}$
	\item $b = 8$
	\item $c = \frac{6}{3}$
	\item $d = 0{,}3333$
	\item $e = \frac{1}{7}$
	\item $f = \frac{1}{8}$
\end{itemize}
\end{exercice}
\begin{exercice}{Affirmations et contre-exemples (Seconde)}
Pour chacune des affirmations suivantes, indiquer si elle est vraie ou fausse. Dans le cas où elle est fausse, donner un contre-exemple.
\begin{itemize}
	\item Si $x \in \N$, alors $2x \in \N$.
	\item Si $x \in \N$, alors $3x + 1 \in \Z$.
	\item Si $x \in \N$, alors $\frac{2x - 5}{2} \in \Z$.
	\item Si $x \in \Z$, alors $\frac{x - 4}{3} \in \D$.
	\item Si $x \in \Q$, alors $x + \frac{1}{2} \in \R$.
	\item Si $x \in \N$, alors $\sqrt{x} \in \Q$.
\end{itemize}
\end{exercice}
\begin{exercice}{Réciproques (Seconde)}
On considère un nombre réel $x$. Pour chacune des affirmations suivantes, indiquer si elle est vraie ou fausse. Indiquer ensuite si sa réciproque est vraie ou fausse.
\begin{itemize}
	\item Si $x \in \N$, alors $x \in \Z$.
	\item Si $x \in \D$, alors $x \in \Q$.
	\item Si $x \in \N$, alors $2x \in \N$.
	\item Si $x \in \Z$, alors $-x \in \N$.
\end{itemize}
\end{exercice}

\end{document}
